\documentclass[11pt]{article}
\usepackage[margin=1in]{geometry}
\usepackage{graphicx}
\usepackage{float}
\usepackage{hyperref}
\usepackage{booktabs}
\usepackage{amsmath}
\usepackage{enumitem}

\title{Breaking Business Boundaries:\\Revenue Forecasting for Vietnamese Fashion E-Commerce}
\author{
  Team GenCore\\[4pt]
  Trinh Hoang Tu (Leader, HUFLIT) \quad
  Nguyen Tan Thang (HCMUS)\\
  Nguyen Trong Huong (UEH) \quad
  Nguyen Minh Nhut (UEH)\\[4pt]
  \texttt{Team Code: YdHJLESH3WgWYCgY8nnb}
}
\date{Datathon 2026 --- VinUniversity}

\begin{document}
\maketitle

\begin{abstract}
We present a zero-leakage forecasting pipeline for predicting daily Revenue and COGS over a 548-day horizon (Jan 2023--Jul 2024) for a Vietnamese fashion e-commerce platform. Our approach combines Seasonal Naive baselines with LightGBM/XGBoost residual correction, calibrated for Tet Nguyen Dan dynamics. On a 548-day holdout, we achieve MAE of 752k VND (Revenue) and 659k VND (COGS). We also provide prescriptive business insights from exploratory analysis of 14 datasets spanning 10 years of operations.
\end{abstract}

%% ════════════════════════════════════════════════════════════════
\section{Part 2: EDA and Prescriptive Analysis}
%% ════════════════════════════════════════════════════════════════

\subsection{Business Objective}

We analyze 14 CSV files (646k orders, 714k line items, 121k customers, 2.4k products) to answer: \textit{where is realized revenue being eroded, and which interventions should be prioritized?}

\subsection{Data Quality Audit}

All 14 tables were audited. Key findings: zero missing values in the sales target; 80\% missing in \texttt{applicable\_category} (promotions table); 12.5\% of orders lack shipment records (cancelled). Foreign key integrity is preserved across all joins. Revenue outliers (169 IQR-flagged days) correspond to legitimate mega-sale events and were retained.

\subsection{Revenue Leakage Drivers}

Three primary leakage drivers were identified:

\begin{enumerate}[leftmargin=*,itemsep=2pt]
  \item \textbf{Cancellation friction (12.5\% of orders):} Credit card payments show the highest cancellation rate. \textit{Action:} implement order confirmation SMS within 30 minutes.
  \item \textbf{Wrong-size returns:} Dominant return reason in Streetwear. \textit{Action:} deploy a size recommendation engine using purchase history.
  \item \textbf{Promotion erosion:} 39\% of items use promotions, but discount depth does not proportionally increase volume. \textit{Action:} cap discount at 25\%, shift budget to retention.
\end{enumerate}

\subsection{Temporal Patterns}

\begin{itemize}[leftmargin=*,itemsep=2pt]
  \item \textbf{Weekly cycle:} Weekdays outperform weekends by 15--20\%.
  \item \textbf{Tet Nguyen Dan:} Revenue surges 3--4 weeks pre-Tet, drops sharply during the holiday week.
  \item \textbf{Mega-sales:} 9.9, 10.10, 11.11, 12.12, and 30/4--1/5 create 2--3x spikes.
  \item \textbf{End-of-month:} Revenue peaks in the last 3--5 days, driven by salary cycles.
\end{itemize}

\subsection{Prescriptive Recommendations (30-60-90 Day Roadmap)}

\begin{itemize}[leftmargin=*,itemsep=2pt]
  \item \textbf{30 days:} Deploy order confirmation SMS; cap promo depth at 25\%.
  \item \textbf{60 days:} Launch size recommendation for Streetwear; pre-position inventory 4 weeks before Tet.
  \item \textbf{90 days:} Build customer segmentation model for targeted retention campaigns.
\end{itemize}

%% ════════════════════════════════════════════════════════════════
\section{Part 3: Forecasting Pipeline and Results}
%% ════════════════════════════════════════════════════════════════

\subsection{Problem Definition}

Forecast daily Revenue and COGS for 548 days (Jan 1, 2023 -- Jul 1, 2024) using historical data from Jul 2012 to Dec 2022. Evaluation metrics: MAE (primary), RMSE, and R$^2$.

\subsection{The Feature Importance Trap}

Our initial XGBoost/LightGBM models used contemporaneous features (web traffic, daily orders, inventory). This achieved CV MAE of 417k but Kaggle MAE of 2M---a 5x degradation. Root cause: these features are unavailable during the 548-day forecast period. Recursive substitution causes exponential error accumulation.

\subsection{Zero-Leakage Architecture}

We redesigned the pipeline with a strict constraint: \textit{every feature must be deterministically computable for any future date}.

\textbf{Tier 1 --- Seasonal Naive Backbone (60\%):}
Blend of 364-day lookback naive (DOW-preserving) and multi-year DOY/DOW median, adjusted by YoY growth rate.

\textbf{Tier 2 --- LightGBM + XGBoost Residual Correction (40\%):}
LightGBM and XGBoost correct naive residuals using only deterministic features: Fourier terms (4 yearly + 2 weekly harmonics), calendar indicators, Tet proximity features, mega-sale flags, historical month/DOW profiles, and auxiliary data as static aggregations. Corrections are damped exponentially ($e^{-h/400}$) at far horizons.

\textbf{Tier 3 --- Tet Calibration:}
Empirical multipliers from 13 years of Tet-relative revenue profiles (day $-30$ to $+20$), applied with optimized strength (Revenue: 0.05, COGS: 0.10).

\subsection{Feature Engineering (Zero-Leakage Verified)}

\begin{itemize}[leftmargin=*,itemsep=1pt]
  \item \textbf{Calendar:} month, DOW, DOY, quarter, weekend, month-start/end
  \item \textbf{Fourier:} $\sin/\cos$ harmonics for yearly and weekly cycles
  \item \textbf{Tet:} days-to-Tet, days-since-Tet, exponential proximity, window flags
  \item \textbf{Mega-sales:} binary flags with $\pm$2 day windows
  \item \textbf{Historical profiles:} month$\times$DOW median, DOY median, rolling tail stats
  \item \textbf{Auxiliary:} web traffic and order count medians by month/DOW (static)
  \item \textbf{Trend:} linear regression projection from historical data
\end{itemize}

\subsection{Validation Strategy}

We use a purged 548-day holdout (Jul 2, 2021 -- Dec 31, 2022) matching the Kaggle test horizon exactly. This eliminates the CV--LB gap from our initial approach. We tested 64 configurations across blend weights, Tet strengths, and training windows.

\subsection{Results}

\begin{table}[H]
\centering
\caption{548-day holdout validation (best of 64 configurations)}
\begin{tabular}{lccc}
\toprule
\textbf{Target} & \textbf{MAE (VND)} & \textbf{RMSE (VND)} & \textbf{R$^2$} \\
\midrule
Revenue & 752,155 & 1,039,000 & 0.552 \\
COGS    & 658,584 &   907,000 & 0.548 \\
\bottomrule
\end{tabular}
\end{table}

Optimal blend: Naive 60\% + Corrected 40\% for both targets. Prophet's direct weight optimized to near zero---its value is as a feature generator, not a standalone forecaster.

\subsection{Feature Importance (SHAP Analysis)}

SHAP analysis of the LightGBM residual corrector reveals the top drivers:

\begin{enumerate}[leftmargin=*,itemsep=2pt]
  \item \textbf{Historical month median} --- seasonal revenue level.
  \item \textbf{Tet effect} --- exponential proximity to Lunar New Year.
  \item \textbf{1-year lag} --- same-period-last-year revenue.
  \item \textbf{DOW median} --- weekday vs.\ weekend differential.
  \item \textbf{Fourier terms} --- annual cyclical patterns.
\end{enumerate}

Full SHAP summary and bar plots are in the GitHub repository (\texttt{output/shap\_*.png}).

\subsection{Error Analysis}

Errors concentrate in: (1) mega-sale spikes (model underestimates peak amplitude); (2) Q2--Q3 normal period (slight over-prediction due to structural differences between holdout and test periods). Distribution shift remains the primary risk.

%% ════════════════════════════════════════════════════════════════
\section{Conclusion}
%% ════════════════════════════════════════════════════════════════

For long-horizon e-commerce forecasting, simple seasonal baselines outperform complex ML models when contemporaneous features are unavailable. Our zero-leakage architecture achieves stable 548-day forecasts by combining Seasonal Naive robustness with gradient boosting pattern correction.

\textbf{Reproducibility:} All code is available at:

\url{https://github.com/thtcsec/Datathon2026-TheGridbreakers-GenCore}

Random seed is fixed at 42. Full pipeline runs in approximately 90 seconds.

\end{document}
